
\emph{Paso 5. Descomposición exacta de $r_{e_0}$.} Por el Paso~3 para $M$ y para $M'$,
$g_M(\theta)-g_{M'}(\theta|_I)=\sum_{e\in F}\phi(t_e)-\theta_r-\theta_c-2$, con
$\theta_r+\theta_c=\log(p_0(1-p_0)/w_{e_0})$, $\sum_F\phi(t_e)=\sum_F\log(1-p_e)+2(2-p_0)$ y
$\sum_{F^*}p_e=2-2p_0$. Por tanto
\begin{equation}\label{eq:accion-hijo}
 H_T(M')-H_T(M)-\log\frac{p_0}{w_{e_0}}=\sum_{e\in F^*}\bigl[-\log(1-p_e)-p_e\bigr]-\mathcal D,\qquad
 \mathcal D=g_{M'}(\theta|_I)-g_{M'}(\theta')\ge0 .
\end{equation}
Para el prefactor, sean $T_{JJ}=\bigl(\begin{smallmatrix}m_r&\tau_0\\ \tau_0&m_c\end{smallmatrix}\bigr)$,
$T_{IJ}$ el bloque de $S_T$, $D_F=\operatorname{diag}(\tau_e)_{a=\lambda(e)}$,
$Q_J=T_{IJ}T_{JJ}^{-1}T_{JI}$ y
\[
 \mathsf S=S'(\tau|)+D_F-Q_J,\qquad W=\mathsf S-S'(\tau') .
\]
Así $\mathsf S$ es el complemento de Schur de $S_T$ en $I$. Para una matriz semidefinida positiva con
núcleo generado por $v$, el menor principal que omite la línea $a$ vale el pseudodeterminante por
$v_a^2/|v|^2$. Aplicado a $S_T$ y a $\mathsf S$, con $|v|^2=2k$ y $|v_I|^2=2k-2$, da
$\mathsf Sv_I=Wv_I=0$ y $\operatorname{pdet}S_T=\tfrac k{k-1}\det T_{JJ}\operatorname{pdet}\mathsf S$.
Resolver $S_T(y,z)=(0,x)$ para $x\perp v_I$ da $\mathsf S^+=\Pi_I^\perp G^0_{II}\Pi_I^\perp$, luego
$\|\mathsf S^+\|_2\le\gamma_A^{-1}$ y $\operatorname{tr}(\mathsf S^+W)=\operatorname{tr}(G^0_{II}W)$.
El término $\tfrac12\log(4\pi k)$ de $L_T$ aporta $\tfrac12\log\tfrac{k-1}k$, que cancela el factor de
tamaño. Con $\ell=L_T-H_T$,
\[
 \ell(M')-\ell(M)=\tfrac12\log\det T_{JJ}+\tfrac12\sum_{e\in F}c_2(p_e)
   -\tfrac12\sum_{e\in M'}[c_2(p'_e)-c_2(p_e)]-\tfrac12\log\det(I-\mathsf S^+W).
\]
Al sumar~\eqref{eq:accion-hijo} obtenemos la identidad exacta
\begin{equation}\label{eq:re-exacto}
 r_{e_0}=-\mathcal D+L_1+L_2+L_3-\log(1+\hat c_{e_0}),
\end{equation}
donde $L_1=\sum_{F^*}[-\log(1-p_e)-p_e]+\tfrac12\log\det T_{JJ}+\tfrac12\sum_Fc_2(p_e)$,
$L_2=-\tfrac12\sum_{M'}[c_2(p'_e)-c_2(p_e)]$ y $L_3=-\tfrac12\log\det(I-\mathsf S^+W)$.

Sea $\beta=G^0a_{e_0}$, de modo que $S_T\beta=a_{e_0}$, $v^\top\beta=0$ y $\mathsf G_0:=\mathsf
G(e_0,e_0)=\beta_r+\beta_c$. Las filas $a\in I$ de $S_T\beta=a_{e_0}$ muestran que
$\rho:=S'(\tau|)\beta_I-b$ es ortogonal a $v_I$, se anula fuera de las líneas $\lambda(e)$, $e\in F^*$,
y
\begin{equation}\label{eq:rho-hijo}
 \rho_{\lambda(e)}=-\varsigma_e-\tau_e(\mathsf G(e,e_0)-1)\quad(e\in F^*),\qquad
 \zeta:=(S'(\tau|)+\Pi')^{-1}\rho+(\xi-\xi_N).
\end{equation}
Como $S'(\tau|)(\xi_N+\beta_I)=\rho$, el vector $\zeta-(\xi+\beta_I)$ es un múltiplo de $v_I$, y
$a_e^\top\zeta=s_e+\mathsf G(e,e_0)$ para toda casilla $e$ de $M'$. Por~\eqref{eq:green-torres}
y~\eqref{eq:newton-hijo}: $|\beta_a|\le2g_G/k$ para $a\in I$, $|\beta_j-1|\le2g_G/k$ para $j\in J$,
$\sum_a|\beta_a|\le2K_G$; $|\mathsf G(e,e_0)-1|\le4g_G/k$ en $F^*$ y $|\mathsf G(e,e_0)|\le4g_G/k$ en
$M'$; y
\begin{equation}\label{eq:cotas-hijo}
 \|\rho\|_\infty\le\tfrac43a^2+\frac{4\bar\tau g_G}k,\qquad
 \|\rho\|_1\le\tfrac83a+\frac{28g_G}{3k},\qquad
 \|\zeta\|_\infty\le\tfrac32K_G\|\rho\|_\infty+61K_G^3a^2 .
\end{equation}

\emph{Paso 6. Cancelación de primer orden.} Sean
\begin{gather*}
 S_1=\sum_{M'}p_e\tau_e\mathsf G(e,e_0),\qquad
 S_2=\sum_{M'}\tau_e\mathsf G(e,e_0)(\mathsf G(e,e)-2),\\
 S_3=\sum_{F^*}\tau_e(\mathsf G(e,e)-2),\qquad S_4=\sum_{F^*}\tau_ep_e .
\end{gather*}
Las identidades siguientes \emph{definen} los restos $E_\bullet$:
\begin{align*}
 \mathcal D&=\tfrac12(m_r\beta_r+m_c\beta_c)-1+\tfrac12\tau_0\mathsf G_0+E_{\mathcal D}, &
 L_1&=\tfrac12(\mu_r+\mu_c)+E_{L1},\\
 L_2&=-S_1+E_{L2}, &
 L_3&=\tfrac12\operatorname{tr}(G^0_{II}W)+E_{L3},\\
 -\log(1+\hat c_{e_0})&=-\hat c_{e_0}+E_{\delta1}, &
 \hat c_{e_0}&=\tfrac12X_{e_0}+E_{\delta2},
\end{align*}
junto con $\operatorname{tr}(G^0_{II}W)=\mathrm{(i)}+\mathrm{(ii)}-\mathrm{(iii)}$, donde
$\mathrm{(i)}=-\sum_{M'}(\tau'_e-\tau_e)\mathsf G(e,e)=-2(1-p_0)+S_2+4S_1+E_{\rm(i)}$,
$\mathrm{(ii)}=\sum_{F^*}\tau_eG^0_{\lambda(e)\lambda(e)}$ y
$\mathrm{(iii)}=\operatorname{tr}(G^0_{II}Q_J)=\tfrac2k-2+m_rG^0_{rr}+m_cG^0_{cc}+2\tau_0G^0_{rc}$, y
con $X_{e_0}=S_3+2S_4+S_2+2S_1-(\mathsf G_0-2)-2p_0+E_X$. La primera identidad es el desarrollo de
Taylor de $g_{M'}$ en $\theta|_I$, es decir, $\mathcal D=-b^\top\xi-\tfrac12\xi^\top S'(\tau|)\xi-R_3$
con $R_3=\sum_{M'}[\phi(t_ee^{s_e})-\phi(t_e)-p_es_e-\tfrac12\tau_es_e^2]$. Sustituyendo
$\xi\equiv\zeta-\beta_I$ y $S'(\tau|)\beta_I=b+\rho$, y usando las filas $r,c$ de
$S_T\beta=a_{e_0}$ en la forma $\sum_{F^*}\tau_e\beta_{\lambda(e)}=2-m_r\beta_r-m_c\beta_c-\tau_0\mathsf
G_0$, resulta
\[
 E_{\mathcal D}=\tfrac12\sum_{F^*}\varsigma_e\beta_{\lambda(e)}-\tfrac12\beta_I^\top\rho+\zeta^\top\rho
   -\tfrac12\zeta^\top S'(\tau|)\zeta-R_3 .
\]
El desarrollo de $\mathrm{(i)}$ usa la conservación exacta de la masa $\sum_{M'}\Delta p_e=1-p_0$,
$\tau=p+\varsigma$ y $s_e=-\mathsf G(e,e_0)+a_e^\top\zeta$; el de $X_{e_0}$
usa~\eqref{eq:chat-torres} y $\varepsilon_e=(\mathsf G(e,e)-2)+2p_e+\tilde\varepsilon_e$. El valor de
$\mathrm{(iii)}$ se sigue de los bloques $(J,I)$ y $(J,J)$ de $S_TG^0=G^0S_T=I-\Pi$.

En la suma de los términos principales de~\eqref{eq:re-exacto}, $S_1$ tiene coeficiente $-1+2-1=0$ y
$S_2$ tiene $\tfrac12-\tfrac12=0$: todo lo que no toca la fila $r$ ni la columna $c$ se cancela. Para
$e\in F^*$ con línea compartida $j$ y $\lambda=\lambda(e)$ se tiene $\mathsf
G(e,e)=G^0_{jj}+G^0_{\lambda\lambda}+2G^0_{j\lambda}$, y la entrada $(r,r)$ de $S_TG^0=I-\Pi$ da
$\sum_{e\in r\setminus e_0}\tau_eG^0_{c(e)r}=1-\tfrac1{2k}-m_rG^0_{rr}-\tau_0G^0_{cr}$, e igualmente para
$c$. Con estas relaciones los términos de orden $1/k$ se cancelan exactamente, y
\begin{align}
 r_{e_0}&=E_{\rm main}-E_{\mathcal D}+E_{L1}+E_{L2}+\tfrac12E_{\rm(i)}+E_{L3}
         +E_{\delta1}-E_{\delta2}-\tfrac12E_X,\label{eq:re-restos}\\
 E_{\rm main}&=\sum_{e\in r\setminus e_0}\frac{p_e^3}{1-2p_e}+\sum_{e\in c\setminus e_0}\frac{p_e^3}{1-2p_e}
   -\tfrac12(G^0_{rr}-1)\mu_r-\tfrac12(G^0_{cc}-1)\mu_c-\tfrac12G^0_{rc}(\mu_r+\mu_c).\notag
\end{align}

\emph{Paso 7. Cotas.} Cada resto de~\eqref{eq:re-restos} es a lo sumo $k^{-2}$ por un polinomio en
$(C_p,K_G,g_G,\gamma_A^{-1})$. Para $E_{\rm main}$ esto se sigue de $\sum_{e\in r}p_e^3/(1-2p_e)\le
\tfrac43a^2$, $|G^0-I|\le g_G/k$ y $\mu_a\le\tfrac43a$. Para $E_{\mathcal D}$ usamos
$\sum_{F^*}\varsigma_e\le\tfrac83a$, las cotas de $\beta$ y~\eqref{eq:cotas-hijo}, $\zeta^\top
S'(\tau|)\zeta=\sum_{M'}\tau_e(a_e^\top\zeta)^2$ con $\sum_e\tau_e\le\tfrac76k$, y
$|\tfrac{d^3}{ds^3}\phi(te^s)|=|\tfrac{d\tau}{ds}|\le1{,}8\tau$, que da $|R_3|\le0{,}35k(6K_Ga)^3$. Para
$E_{L1}$: $|c_2(u)+u^2|\le3{,}6u^3$, $0\le-\log(1-u)-u-\tfrac12u^2\le0{,}39u^3$,
$|\log m_a-\mu_a|\le\tfrac12\mu_a^2$ y $\tau_0^2\le(\tfrac76a)^2$. Para $E_{L2}$: $|\dot c_2(u)/2+u|\le
4{,}2u^2$, $|\Delta p_e-\tau_es_e|\le0{,}9\tau_es_e^2$ y $|\ddot c_2|\le5{,}3$. Para $E_{\rm(i)}$:
$\Delta\varsigma=\dot\varsigma\Delta p+\tfrac12\ddot\varsigma\Delta p^2$ y $|\dot\varsigma(p)\mathsf
G(e,e)-4p|\le22{,}5p^2+8pg_G/k$. Para $E_{L3}$: con $Y=(\mathsf S^+)^{1/2}W(\mathsf S^+)^{1/2}$ se tiene
$\|Y\|\le\|W\|_\infty/\gamma_A\le\tfrac12$, $\|Y\|_F\le\|W\|_F/\gamma_A$ y
$|\log\det(I-Y)+\operatorname{tr}Y|\le\|Y\|_F^2$. Por la conservación de la masa en cada línea $a\in I$,
las entradas de $W$ son $W_{aa}=\varsigma_{e_a}-\sum_{e\in a\cap M'}\Delta\varsigma_e-(Q_J)_{aa}$, con
$e_a$ la casilla de $F^*$ en $a$, y $W_{ab}=-(\tau'_e-\tau_e)-(Q_J)_{ab}$ para la casilla $e$ de las
líneas $a,b$; por tanto $\max|W_{ab}|=O(k^{-2})$, $\|W\|_\infty=O(k^{-1})$ y $\|W\|_F^2=O(k^{-2})$. Por
último, $|E_{\delta1}|\le\hat c_{e_0}^2$, $|E_{\delta2}|\le\tfrac23a|X_{e_0}|$ porque
$\tau_0/p_0-1\le\tfrac43a$, y $E_X$ reúne $\tau_0\varepsilon_0\mathsf G_0$, las diferencias $\mathsf
G(e,e_0)-1$ en $F^*$ y los términos $\tilde\varepsilon$, todos controlados por el Paso~3. Sumar las
cotas en $k=k_1$ y usar que cada una es $k^{-2}$ o $k^{-3}$ por una constante da $|r_{e_0}|\le C_r/k^2$
para todo $k\ge k_1$. El enunciado para $\mathcal K$ se sigue de las propiedades de la clase
estacionaria recordadas en la Sección~\ref{sec:torres}.
\end{proof}
